\chapter*{Fontes}
\addcontentsline{toc}{chapter}{Fontes}
\markboth{Fontes}{Fontes}
\small
Os números sobrescritos no texto remetem a esta lista. As questões trazem a procedência no próprio cabeçalho. Os endereços da web foram acessados em novembro de 2026, e em todas as páginas da web (Wikipédia, enciclopédias e repositórios) vale a versão dessa data. Quando a fonte acessível é secundária, a obra original vem citada ao lado, sempre que localizada.

\begingroup
\makeatletter
% a lista entra logo abaixo do título "Fontes": sem novo capítulo nem cabeçalho "Referências"
\renewcommand{\chapter}{\@ifstar{\@gobble}{\@gobble}}
\renewcommand{\@mkboth}[2]{}
\makeatother
\begin{thebibliography}{999}\setlength{\itemsep}{0.2em}
\bibitem{platoSoph} Platão, \emph{Sofista}, 262--263 (trad. H.~N.~Fowler, 1921). \url{https://cts.perseids.org/read/greekLit/tlg0059/tlg007/perseus-eng2/263}
\bibitem{sepAncient} S.~Bobzien, \enquote{Ancient Logic}, \emph{Stanford Encyclopedia of Philosophy}. \url{https://plato.stanford.edu/archives/Win2012/entries/logic-ancient/}
\bibitem{aristInt} Aristóteles, \emph{Da Interpretação}, cap.~4 (trad. E.~M.~Edghill). \url{https://www.ellopos.net/elpenor/greek-texts/ancient-greece/aristotle/interpretation.asp?pg=4}
\bibitem{aristInt7} Aristóteles, \emph{Da Interpretação}, cap.~7, 17b (trad. E.~M.~Edghill), Elpenor. \url{https://www.ellopos.net/elpenor/greek-texts/ancient-greece/aristotle/interpretation.asp?pg=16}
\bibitem{sepContr} L.~Horn, \enquote{Contradiction}, \emph{Stanford Encyclopedia of Philosophy}. \url{https://plato.stanford.edu/archives/fall2008/entries/contradiction/}
\bibitem{wedin} M.~V.~Wedin, \enquote{Aristotle on the Firmness of the Principle of Non-Contradiction}, \emph{Phronesis}. \url{https://sites.unimi.it/zucchi/NuoviFile/Wedin04.pdf}
\bibitem{futcont} \enquote{Problem of future contingents}, Wikipedia. \url{https://en.wikipedia.org/wiki/Problem_of_future_contingents}
\bibitem{sepFut} S.~Knuuttila, \enquote{Medieval Theories of Future Contingents}, \emph{Stanford Encyclopedia of Philosophy}. \url{https://plato.stanford.edu/archives/spr2024/entries/medieval-futcont/}
\bibitem{megarians} \enquote{Megarians}, \emph{Encyclopedia of Philosophy} (Encyclopedia.com). \url{https://www.encyclopedia.com/humanities/encyclopedias-almanacs-transcripts-and-maps/megarians}
\bibitem{oupEub} \enquote{Eubulides and his paradoxes}, OUPblog, 2017. \url{https://blog.oup.com/2017/08/eubulides-paradoxes-philosophy/}
\bibitem{bobzien} S.~Bobzien, \enquote{Stoic Logic}, in \emph{The Cambridge Companion to the Stoics}, 2003. \url{https://philarchive.org/archive/BOBSL}
\bibitem{stoicNeg} S.~Lombardi, \enquote{La negación externa estoica como negación existencial}, \emph{Anales del Seminario de Historia de la Filosofía} 43(2), 2026. \url{https://revistas.ucm.es/index.php/ASHF/article/view/99170}
\bibitem{dlDouble} Diógenes Laércio, \emph{Vidas}, VII, 69 (a dupla negação estoica; trad. R.~D.~Hicks, 1925). \url{https://fgh.perseids.org/read/greekLit/tlg0004/tlg001/perseus-eng2/7.1.61-7.1.80}
\bibitem{beziau} J.-Y.~Beziau, \enquote{Is the principle of contradiction a consequence of $x^2=x$?}, \emph{Logica Universalis}, 2018. \url{https://core.ac.uk/download/148070305.pdf}
\bibitem{maaBoole} Resenha de G.~Boole, \emph{The Laws of Thought}, MAA Reviews. \url{https://www.maa.org/node/106414}
\bibitem{miller} J.~Miller, \emph{Earliest Uses of Symbols of Set Theory and Logic} (MacTutor). \url{https://mathshistory.st-andrews.ac.uk/Miller/mathsym/set/}
\bibitem{sepBrouwer} M.~van Atten, \enquote{Luitzen Egbertus Jan Brouwer}, \emph{Stanford Encyclopedia of Philosophy}. \url{https://plato.stanford.edu/entries/brouwer/}
\bibitem{moscho} J.~R.~Moschovakis; G.~Vafeiadou, \enquote{Intuitionistic Mathematics and Logic}, arXiv:2003.01935. \url{https://arxiv.org/pdf/2003.01935}
\bibitem{wolenski} J.~Woleński, \enquote{Logical Ideas of Jan Łukasiewicz}, \emph{Studia Humana} 8(2), 2019, p.~3--7. \url{https://www.studiahumana.com/pliki/wydania/2%20Wolenski.pdf}
\bibitem{openlogic} Open Logic Project, \enquote{Łukasiewicz logic}. \url{https://builds.openlogicproject.org/locale/pt/content/many-valued-logic/three-valued-logics/lukasiewicz.pdf}
\bibitem{educapes} M.~R. Kasper, \emph{Proposta didática: utilizando questões da OBMEP para o ensino de lógica formal}, produto de mestrado (PROFMAT), IFRS -- Campus Canoas, 2025 (EduCAPES). \url{https://educapes.capes.gov.br/bitstream/capes/1131738/2/Produto-Marcelo.pdf}
\bibitem{matos} G.~Matos, \enquote{Negative Concord and the Scope of Negation}, \emph{Catalan Working Papers in Linguistics} 7, 1999. \url{https://core.ac.uk/download/13268562.pdf}
\bibitem{correio} \enquote{Por que três palavras negativas não viram afirmação}, \emph{Correio Braziliense} (com base no Ciberdúvidas). \url{https://www.correiobraziliense.com.br/cbradar/por-que-tres-palavras-negativas-nao-viram-afirmacao-gramatica/}
\bibitem{goldbach} T.~Oliveira e Silva; S.~Herzog; S.~Pardi, \enquote{Empirical verification of the even Goldbach conjecture and computation of prime gaps up to $4\cdot10^{18}$}, \emph{Mathematics of Computation} 83(288), 2014, p.~2033--2060. \url{https://www.ams.org/journals/mcom/2014-83-288/S0025-5718-2013-02787-1}
\bibitem{smullyan} \enquote{Knights and Knaves} (R.~Smullyan, \emph{What Is the Name of This Book?}, 1978), Wikipedia. \url{https://en.wikipedia.org/wiki/Knights_and_Knaves}
\bibitem{smullyan1978} R.~Smullyan, \emph{What Is the Name of This Book? The Riddle of Dracula and Other Logical Puzzles}. Englewood Cliffs: Prentice-Hall, 1978.
\bibitem{obmep2010} OBMEP 2010, 1ª fase — banco de questões de raciocínio lógico (Olimpédia). \url{https://olimpedia.fandom.com/pt-br/wiki/Banco_de_Quest%C3%B5es_-_Racioc%C3%ADnio_L%C3%B3gico_-_N%C3%ADvel_2}
\bibitem{britIncl} \enquote{Inclusive disjunction}, \emph{Encyclopaedia Britannica}. \url{https://www.britannica.com/topic/inclusive-disjunction}
\bibitem{plutarchNum} E.~W.~Weisstein, \enquote{Plutarch Numbers}, \emph{MathWorld}. \url{https://mathworld.wolfram.com/PlutarchNumbers.html}
\bibitem{mathpages} \enquote{Hipparchus on Compound Statements}, \emph{MathPages}. \url{https://mathpages.com/home/kmath397/kmath397.htm}
\bibitem{bobzienComb} S.~Bobzien, \enquote{The Combinatorics of Stoic Conjunction: Hipparchus Refuted, Chrysippus Vindicated}, \emph{Oxford Studies in Ancient Philosophy} 40, 2011. \url{https://philarchive.org/archive/BOBTCO-2}
\bibitem{britBoole} \enquote{History of logic: Boole and De Morgan}, \emph{Encyclopaedia Britannica}. \url{https://www.britannica.com/topic/history-of-logic/Boole-and-De-Morgan}
\bibitem{ieeeJevons} \enquote{W.~Stanley Jevons}, \emph{Computer Pioneers}, IEEE Computer Society. \url{https://history.computer.org/pioneers/jevons.html}
\bibitem{waszek} D.~Waszek, edição comentada do manuscrito de Boole \enquote{On the Nature of Thought}, \emph{History and Philosophy of Logic}, 2025. \url{https://www.normalesup.org/~waszek/docs/Waszek-Boole-NoT.pdf}
\bibitem{couturat} L.~Couturat, \emph{The Algebra of Logic} (trad. inglesa), Project Euclid. \url{https://projecteuclid.org/download/pdf_1/euclid.chmm/1263313795}
\bibitem{jevons1870} W.~S.~Jevons, \enquote{On the Mechanical Performance of Logical Inference}, \emph{Philosophical Transactions of the Royal Society} 160, 1870, p.~497--518. DOI 10.1098/rstl.1870.0022.
\bibitem{macJevons} \enquote{William Stanley Jevons}, MacTutor History of Mathematics. \url{https://mathshistory.st-andrews.ac.uk/Biographies/Jevons/}
\bibitem{venn} J.~Venn, \enquote{On the Diagrammatic and Mechanical Representation of Propositions and Reasonings}, \emph{Philosophical Magazine}, série 5, 10(59), 1880, p.~1--18. \url{https://doi.org/10.1080/14786448008626877}
\bibitem{vennSurvey} F.~Ruskey; M.~Weston, \enquote{Venn Diagram Survey: John Venn}, \emph{Electronic Journal of Combinatorics}, DS5. \url{https://www.combinatorics.org/files/Surveys/ds5/VennJohnEJC.html}
\bibitem{vennWiki} \enquote{Venn diagram} (seção \emph{History}), Wikipedia. \url{https://en.wikipedia.org/wiki/Venn_diagram}
\bibitem{fregeColor} M.~Potter, \enquote{Propositions in Wittgenstein and Ramsey}, Univ. de Cambridge (repositório Apollo/CORE); trecho sobre o \enquote{colorido} (\emph{Färbung}) de \enquote{e} e \enquote{mas} em Frege. \url{https://core.ac.uk/download/237398527.pdf}
\bibitem{shannonMIT} C.~E.~Shannon, \emph{A Symbolic Analysis of Relay and Switching Circuits}, dissertação de mestrado, MIT, 1937. \url{https://dspace.mit.edu/handle/1721.1/11173}
\bibitem{liversidge} M.~Dawson, \emph{Mind, Body, World}, seção 2.4 (cita entrevista de Shannon a A.~Liversidge, 1993), LibreTexts. \url{https://socialsci.libretexts.org/Bookshelves/Psychology/Cognitive_Psychology/Mind_Body_World_-_Foundations_of_Cognitive_Science_(Dawson)/02\%3A_Multiple_Levels_of_Investigation/2.04\%3A_From_the_Formal_to_the_Physical}
\bibitem{shannonConv} M.~Bullynck, \enquote{Switching the Engineer's Mind-Set to Boolean}, in T.~Haigh (org.), \emph{Exploring the Early Digital}, Springer, 2019. \url{https://link.springer.com/content/pdf/10.1007/978-3-030-02152-8_5.pdf}
\bibitem{elkindZach} L.~D.~C.~Elkind; R.~Zach, \enquote{The Genealogy of `$\vee$'}, \emph{The Review of Symbolic Logic} 16(3), 2023. \url{https://arxiv.org/abs/2012.06072}
\bibitem{obmep2009} OBMEP 2009, 1ª fase, nível 3 (prova). \url{http://obmep2010.obmep.org.br/provas_static/export/sites/default/provas/provas2009/1_fase/prova_1fase_n3_final.pdf}
\bibitem{obmep2005} OBMEP 2005, 1ª fase, nível 3, questão 20 (prova; cópia hospedada pelo IFES). \url{https://st.ifes.edu.br/images/stories/Not\%C3\%ADcias/2021/obmep/prova_nivel3_32033486.pdf}
\bibitem{sepDialectical} S.~Bobzien, \enquote{Dialectical School}, \emph{Stanford Encyclopedia of Philosophy}, 2004 (ed. verão 2021). \url{https://plato.stanford.edu/archives/sum2021/entries/dialectical-school/}
\bibitem{ls35b} Sexto Empírico, \emph{Hipotiposes Pirrônicas} 2.110--113 (texto LS 35B, trad. Long e Sedley), material de curso da Universidade de Atenas. \url{https://eclass-b.uoa.gr/modules/document/file.php/PHS582/Stoic\%20Epistemology/plus\%20texts.proof.docx}
\bibitem{whiteCrows} M.~J.~White, \enquote{What Worried the Crows?}, \emph{The Classical Quarterly}. \url{https://www.cambridge.org/core/journals/classical-quarterly/article/what-worried-the-crows/B202EADB2A9589F6F673E71FF5DE9D88}
\bibitem{organonF} M.~Cresswell, \enquote{Modal Logic before Kripke}, \emph{Organon F} 26(3), 2019, p.~323--339. \url{https://www.sav.sk/journals/uploads/08241202orgf.2019.26302.pdf}
\bibitem{openlogicPar} Open Logic Project, \enquote{Paradoxes of the Material Conditional}. \url{https://builds.openlogicproject.org/locale/pt/content/counterfactuals/introduction/paradoxes-material.pdf}
\bibitem{wasonWiki} \enquote{Wason selection task}, Wikipedia. \url{https://en.wikipedia.org/wiki/Wason_selection_task}
\bibitem{wasonShapiro} P.~C. Wason; D. Shapiro, \enquote{Natural and contrived experience in a reasoning problem}, \emph{Quarterly Journal of Experimental Psychology} 23(1), 1971, p.~63--71 (Exp.~II, Tabela VI). DOI 10.1080/00335557143000068.
\bibitem{wason1968} P.~C.~Wason, \enquote{Reasoning about a rule}, \emph{Quarterly Journal of Experimental Psychology} 20(3), 1968, p.~273--281. DOI 10.1080/14640746808400161.
\bibitem{griggsCox1982} R.~A.~Griggs; J.~R.~Cox, \enquote{The elusive thematic-materials effect in Wason's selection task}, \emph{British Journal of Psychology} 73(3), 1982, p.~407--420. DOI 10.1111/j.2044-8295.1982.tb01823.x.
\bibitem{sussexWason} \enquote{Reasoning in the Wason 4-Card Selection Task}, material de aula, Universidade de Sussex. \url{https://users.sussex.ac.uk/~grahamh/RM1web/Wason\%20materials2010.pdf}
\bibitem{cosmides1989} L.~Cosmides, \enquote{The logic of social exchange: Has natural selection shaped how humans reason? Studies with the Wason selection task}, \emph{Cognition} 31(3), 1989, p.~187--276. DOI 10.1016/0010-0277(89)90023-1.
\bibitem{sperber} D.~Sperber; V.~Girotto, \enquote{Does the selection task detect cheater detection?} (versão disponível no repositório CORE). \url{https://core.ac.uk/download/46683914.pdf}
\bibitem{ravenWiki} \enquote{Raven paradox}, Wikipedia. \url{https://en.wikipedia.org/wiki/Raven_paradox}
\bibitem{hempel1945} C.~G.~Hempel, \enquote{Studies in the Logic of Confirmation (I.)}, \emph{Mind} 54(213), 1945, p.~1--26. DOI 10.1093/mind/LIV.213.1.
\bibitem{fitelson} B.~Fitelson; J.~Hawthorne, \enquote{How Bayesian Confirmation Theory Handles the Paradox of the Ravens}, in E.~Eells; J.~Fetzer (orgs.), \emph{The Place of Probability in Science}, Springer, 2010 (versão em manuscrito). \url{https://mail.joelvelasco.net/teaching/5311/FitelsonHawthorneRavens.pdf}
\bibitem{obmep2007} OBMEP 2007, 1ª fase, nível 2 (prova). \url{http://obmep2010.obmep.org.br/provas_static/export/sites/default/provas/provas2007/2007/1_fase/N2_curvas_final_imp_2507.pdf}
\bibitem{simulado2017N3} Simulado OBMEP 2017, nível 3 (soluções). \url{https://www.obmep.org.br/docs/sim-2017-N3-solucoes.pdf}
\bibitem{truthTableWiki} \enquote{Truth table}, Wikipedia (seção de história, com a pesquisa de I.~Anellis e J.~Shosky). \url{https://en.wikipedia.org/wiki/Truth_table}
\bibitem{anellis} I.~H.~Anellis, \enquote{Peirce's Truth-Functional Analysis and the Origin of Truth Tables}, arXiv:1108.2429. \url{https://arxiv.org/abs/1108.2429}
\bibitem{wittND} J.~Speaks, notas de aula sobre o \emph{Tractatus} (proposições elementares; tautologias e contradições, 4.46--4.4661), Univ. de Notre Dame. \url{https://www3.nd.edu/~jspeaks/courses/2007-8/43904/_HANDOUTS/wittgenstein-elementary-propositions.pdf}
\bibitem{wittTract} L.~Wittgenstein, \emph{Tractatus Logico-Philosophicus}, trad. C.~K.~Ogden, 1922 (domínio público), prop.~4.46--4.461. \url{https://www.gutenberg.org/ebooks/5740}
\bibitem{sparkTract} \enquote{Tractatus Logico-philosophicus: Section 7}, SparkNotes. \url{https://www.sparknotes.com/philosophy/tractatus/section7/}
\bibitem{booleAlgWiki} \enquote{Boolean algebra}, Wikipedia. \url{https://en.wikipedia.org/wiki/Boolean_algebra}
\bibitem{shefferWiki} \enquote{Henry M. Sheffer}, Wikipedia. \url{https://en.wikipedia.org/wiki/Henry_M._Sheffer}
\bibitem{sheffer1913} H.~M.~Sheffer, \enquote{A set of five independent postulates for Boolean algebras, with application to logical constants}, \emph{Transactions of the American Mathematical Society} 14(4), 1913, p.~481--488.
\bibitem{nicod1917} J.~Nicod, \enquote{A reduction in the number of primitive propositions of logic}, \emph{Proceedings of the Cambridge Philosophical Society} 19, 1917, p.~32--41.
\bibitem{shefferStroke} \enquote{Sheffer stroke} (seção \emph{History}), Wikipedia. \url{https://en.wikipedia.org/wiki/Sheffer_stroke}
\bibitem{carroll1896} L.~Carroll, \emph{Symbolic Logic}, parte I, 1896 (domínio público; Project Gutenberg \#28696). \url{https://www.gutenberg.org/ebooks/28696}
\bibitem{cook1971} S.~A.~Cook, \enquote{The complexity of theorem-proving procedures}, \emph{Proceedings of the 3rd Annual ACM Symposium on Theory of Computing (STOC)}, 1971, p.~151--158. DOI 10.1145/800157.805047.
\bibitem{satWiki} \enquote{Boolean satisfiability problem}, Wikipedia. \url{https://en.wikipedia.org/wiki/Boolean_satisfiability_problem}
\bibitem{fatec2013} FATEC, vestibular do 1º semestre de 2013, questão de lógica (transcrição: Estratégia Vestibulares). \url{https://vestibulares.estrategia.com/public/questoes/Logica-tem-que36f11d49c9/}
\bibitem{sepFOL} W.~Ewald, \enquote{The Emergence of First-Order Logic}, \emph{Stanford Encyclopedia of Philosophy}. \url{https://plato.stanford.edu/entries/logic-firstorder-emergence/}
\bibitem{euler1732} L.~Euler, \enquote{Observações sobre um teorema de Fermat e outros relativos aos números primos} (1732; E26), tradução inglesa, arXiv:math/0501118. \url{https://arxiv.org/abs/math/0501118}
\bibitem{millersville} B.~Ikenaga, \enquote{Fermat Numbers}, notas de teoria dos números, Millersville University. \url{https://sites.millersville.edu/bikenaga/number-theory/fermat-numbers/fermat-numbers.html}
\bibitem{perrinEuler} D.~Perrin, \enquote{Pourquoi y a-t-il beaucoup de nombres premiers de la forme $n^2+n+41$?}, Université Paris-Saclay. \url{https://www.imo.universite-paris-saclay.fr/~perrin/journeedu2311/redaction2311e.pdf}
\bibitem{sepSquare} T.~Parsons, \enquote{The Traditional Square of Opposition}, \emph{Stanford Encyclopedia of Philosophy}, 1997 (ed. primavera 2008). \url{https://plato.stanford.edu/archives/spr2008/entries/square/}
\bibitem{deMorganWiki} \enquote{De Morgan's laws}, Wikipedia (seções de história, engenharia e busca de texto). \url{https://en.wikipedia.org/wiki/De_Morgan\%27s_laws}
\bibitem{aristPriorAn} Aristóteles, \emph{Primeiros Analíticos} I.23, 41a23--30 (trad. inglesa de A.~J.~Jenkinson, domínio público), Elpenor. \url{https://www.ellopos.net/elpenor/greek-texts/ancient-greece/aristotle/prior-analytics.asp?pg=41}
\bibitem{corry} L.~Corry, nota sobre Aristóteles e a incomensurabilidade da diagonal, Universidade de Tel Aviv. \url{https://www.tau.ac.il/~corry/Narrative/notes/aristo.html}
\bibitem{fatec2019} FATEC, vestibular do 1º semestre de 2019, questão 11 (resolução comentada do Curso Objetivo). \url{https://www.curso-objetivo.br/vestibular/resolucao_comentada/fatec/2019/fatec2019.pdf}
\bibitem{fatec2015} FATEC, vestibular do 1º semestre de 2015 (prova de dez./2014), questão 11 (resolução comentada do Curso Objetivo). \url{https://curso-objetivo.br/vestibular/resolucao_comentada/fatec/2015/fatec2015.pdf}
\bibitem{fatec2014s2} FATEC, vestibular do 2º semestre de 2014, questão 12 (resolução comentada do Curso Objetivo). \url{https://curso-objetivo.br/vestibular/resolucao_comentada/fatec/2014_2/fatec2014_2.pdf}
\bibitem{fme} G.~Iezzi; C.~Murakami, \emph{Fundamentos de Matemática Elementar}, vol.~1: \emph{Conjuntos, funções}, cap.~I, \enquote{Noções de lógica}. São Paulo: Atual.
\end{thebibliography}
\endgroup
